\begin{frame}{정밀도 축을 작업별로 대입하면}
  \small
  \begin{tabular}{@{}p{4.6cm}p{2.8cm}p{6.6cm}@{}}
    \toprule
    작업 (Week) & 필요한 정밀도 & 이유 \\
    \midrule
    CartPole, Pendulum (9$\sim$11) & 낮음$\sim$중간 &
      단일 막대/진자. 접촉 없음, 물리 근사만으로도 충분 \\
    Reacher, 2관절 로봇팔 (13) & 중간 & 관절 운동은 있지만
      \textbf{자발적인 강한 접촉} 없음 \\
    사족보행/인간형 보행 (13.2) & \textbf{높음} & 발이 바닥에
      \textbf{붙었다 떨어지는} 반복 접촉. 접촉력 정확도가 핵심 \\
    카메라로 판단하는 정책 & 높음(센서) + 물리 & 렌더링
      리얼리즘과 물리 정밀도 \textbf{둘 다} 필요 \\
    \bottomrule
  \end{tabular}
  \vskip0.7em
  \begin{block}{정밀도 축의 결론}
    ``\textbf{접촉이 얼마나 중요한가}''로 결정된다. 접촉이 거의
    없으면(Pendulum) Gymnasium 기본 환경, 관절이 있지만 강한
    접촉이 없으면(Reacher) PyBullet/MuJoCo 어느 쪽이든 충분,
    \textbf{발이 바닥을 차는 보행}이라면 MuJoCo급 정밀 LCP
    접촉이 \textbf{필수}.
  \end{block}
\end{frame}
